% Actual class: 09
% Slide deck: Part 4 (Process Synchronization)
% Exact scope: pp. 28--36
% Video supplement: Part 4.5 and Part 4.6 complete
% Human verified: Yes

\section{第 09 节课：Semaphore 类型与实现}
\label{lecture:SC2005-L09}

\lecturemeta
  {SC2005-L09}
  {2026-09-10}
  {Part 4 (Process Synchronization)}
  {主讲义 pp.~28--36}
  {Part 4.5 Types of Semaphores 与 Part 4.6 Semaphore Implementations 完整}
  {已确认（2026-09-10）}

\subsection{本讲主线}

上一讲的 TestAndSet（测试并置位）提供底层 atomic read--modify--write（原子读改写），但程序仍需自行管理 lock 与等待方式。本讲引入 semaphore（信号量）：只允许通过 atomic \texttt{Wait} 与 \texttt{Signal} 访问的同步对象。Binary semaphore 用于 mutual exclusion；counting semaphore 管理多个同类资源。等待既可用 spinning 实现，也可通过 OS queue 将 process 置为 waiting；后者不浪费长时间的 CPU cycles，但会产生 context-switch overhead。

\lecturetopic{SC2005-L09-T01}{Semaphore 抽象与两种类型}

\keyidea{Semaphore 不是可任意读写的普通整数，而是“状态 $+$ 两个 atomic operations”的同步抽象。}

课程使用如下抽象语义：
\[
  \texttt{Wait(S)}:
  \begin{cases}
    S\gets S-1, & S>0,\\
    \text{wait}, & S\le 0,
  \end{cases}
  \qquad
  \texttt{Signal(S)}:S\gets S+1.
\]
\texttt{Wait} 也称 P/down，用于 acquire permit（获取许可）；\texttt{Signal} 也称 V/up，用于 release permit（释放许可）。对 $S$ 的检查、修改以及必要的等待或唤醒必须构成不可分割的 protocol。

\begin{itemize}
  \item \textbf{Binary semaphore（二元信号量）：}逻辑上的可用许可最多为 $1$，常用于 mutual exclusion。
  \item \textbf{Counting semaphore（计数信号量）：}初值通常等于某类 shared resource 的实例数，每次 \texttt{Wait} 消耗一个实例，每次 \texttt{Signal} 归还一个实例。
\end{itemize}

Blocking implementation 中的内部 $S.value$ 可以为负；这不表示“负的资源”，而表示尚有等待请求。若 $S.value=-k$，则等待队列中有 $k$ 个 processes。

\lecturetopic{SC2005-L09-T02}{Binary Semaphore 实现 Mutual Exclusion}

共享 semaphore \texttt{mutex} 初始化为 $1$：
\[
  \texttt{Wait(mutex)}\rightarrow
  \text{critical section}\rightarrow
  \texttt{Signal(mutex)}\rightarrow
  \text{remainder section}.
\]
第一个 process 执行 \texttt{Wait} 后取得唯一许可；其他 process 必须等待 holder 执行 \texttt{Signal}。\texttt{Signal} 只会把某个等待者从 waiting 变为 ready，不保证它立即成为 running。Holder 离开 critical section 后继续执行 remainder section；释放 semaphore 不等于结束整个 process。

\textbf{对应独立检验例：}\relatedexercise{SC2005-L09-E01}{Binary Semaphore 的状态轨迹}。

\lecturetopic{SC2005-L09-T03}{Busy-Waiting Semaphore 与 Atomicity}

Classical busy-waiting implementation 为
\[
\begin{aligned}
  \texttt{Wait(S):}\;&\texttt{while (S <= 0);}\quad \texttt{S--;}\\
  \texttt{Signal(S):}\;&\texttt{S++;}
\end{aligned}
\]
等待者不断读取 $S$，属于 spinning（自旋）。它避免 block/wakeup 引起的 context switch，适合 multicore 上预期极短的等待；若 critical section 较长，则会浪费 CPU cycles。

仅保证一次 \texttt{S--} atomic 并不够；“检查 $S>0$”与“递减 $S$”也必须合成一个 atomic operation。否则两个 processes 可先后通过同一次 $S=1$ 的检查，再分别执行 decrement，最终同时进入 critical section。

\textbf{对应独立检验例：}\relatedexercise{SC2005-L09-E02}{检查与递减分离造成的 Mutual-Exclusion Violation}。

\lecturetopic{SC2005-L09-T04}{Blocking Semaphore：Value 与 Waiting Queue}

Blocking implementation 将 semaphore 表示为
\[
  S=(\texttt{value},L),
\]
其中 $L$ 保存等待该 semaphore 的 PCB list。\texttt{block()} 将当前 process 从 ready queue 移入 $L$ 并置为 waiting；\texttt{wakeup()} 从 $L$ 取出一个 process，放回 ready queue 并置为 ready。

\[
\begin{aligned}
  \texttt{Wait(S):}\quad
    &\texttt{S.value--;}\\[-0.2em]
    &\texttt{if (S.value < 0) block();}\\[0.35em]
  \texttt{Signal(S):}\quad
    &\texttt{S.value++;}\\[-0.2em]
    &\texttt{if (S.value <= 0) wakeup(P);}
\end{aligned}
\]

\begin{figure}[htbp]
  \centering
  \resizebox{0.98\textwidth}{!}{\begin{tikzpicture}[
  >=Latex,
  state/.style={draw=noteBlue,very thick,rounded corners=3pt,fill=blue!5,text width=2.55cm,minimum height=2.05cm,align=center,font=\small,inner sep=5pt},
  blocked/.style={state,draw=ntuRed,fill=red!6},
  ready/.style={state,draw=green!45!black,fill=green!7},
  flow/.style={-{Latex[length=2.2mm]},thick,draw=black!65}
]
  \node[state] (s0) at (0,0)
    {Initial\\$S.value=1$\\$L=[\ ]$};
  \node[state] (s1) at (3.4,0)
    {$P_0$: Wait\\$S.value=0$\\$P_0$: CS};
  \node[blocked] (s2) at (6.8,0)
    {$P_1$: Wait\\$S.value=-1$\\$L=[P_1]$};
  \node[blocked] (s3) at (10.2,0)
    {$P_2$: Wait\\$S.value=-2$\\$L=[P_1,P_2]$};
  \node[ready] (s4) at (13.6,0)
    {$P_0$: Signal\\$S.value=-1$\\$P_1$: ready, $L=[P_2]$};

  \draw[flow] (s0) -- (s1);
  \draw[flow] (s1) -- (s2);
  \draw[flow] (s2) -- (s3);
  \draw[flow] (s3) -- (s4);
\end{tikzpicture}
}
  \caption{初值 $S.value=1$ 的 blocking semaphore 状态轨迹。图中等待队列假定采用 FIFO；wakeup 后 $P_1$ 仅进入 ready。}
  \label{fig:sc2005-l09-blocking-semaphore}
\end{figure}

先 decrement 再检查使 $S.value$ 同时记录当前申请：非负值表示剩余 permits，负值的绝对值表示 queued waiters。执行 \texttt{Signal} 后若结果仍不大于 $0$，说明至少有一个请求尚未满足，需要唤醒一个等待者。

\lecturetopic{SC2005-L09-T05}{Blocking 实现的原子边界与选择}

Blocking 并未取消 atomicity 要求。\texttt{Wait} 中的 decrement、test 与 enqueue/block protocol，以及 \texttt{Signal} 中的 increment、test 与 dequeue/wakeup protocol，必须相对于彼此保持原子性；否则会发生 lost wakeup、progress failure，甚至 mutual-exclusion violation。

不能先 \texttt{block()}、再执行 \texttt{S.value--}：在两者之间，另一个 process 可能 \texttt{Signal} 并唤醒当前 process，第三个 process 又取得刚释放的许可；当前 process 恢复后若继续 decrement 并进入，便会与第三个 process 同时持有同一许可。

实现选择取决于 expected waiting time，而不只是 CPU core 数：
\begin{center}
\small
\begin{tabularx}{0.96\textwidth}{@{}>{\raggedright\arraybackslash}p{2.55cm}XX@{}}
  \toprule
  Strategy & Cost & 适用情形 \\
  \midrule
  Busy waiting / spinning & 持续占用一个 CPU core；无 block/wakeup context switch & Multicore，holder 可在另一 core 运行，且等待预计极短 \\
  Blocking & 需要 scheduler、queue 与 context switch & Single-core、临界区较长或等待时间不可预测 \\
  \bottomrule
\end{tabularx}
\end{center}

实际 blocking semaphore 的内部短 critical section 仍可由 hardware atomic instruction 或短暂 spinlock 保护；被消除的是等待资源期间的长时间 spinning，而不是一切底层原子机制。

\textbf{对应独立检验例：}\relatedexercise{SC2005-L09-E03}{Spinning 与 Blocking 的选择}。

\subsection{一分钟回顾}

Semaphore 用 atomic \texttt{Wait}/\texttt{Signal} 管理进入许可：binary semaphore 保护单一 critical section，counting semaphore 管理多个同类资源。Busy waiting 避免 context switch 但消耗 CPU；blocking implementation 用 $S.value$ 与 waiting queue 记录未满足请求。无论采用哪种等待方式，对 semaphore 状态的检查、更新和队列操作都必须保持原子性；\texttt{wakeup} 只使 process 进入 ready，而非立即 running。
